% Limit: for every band d +- eta there is a strip -delta < h < delta (h != 0) where the graph stays in the band.
% g(h) = 2 + 0.5 h + 0.25 h^2, eta = 0.5, delta = 0.7 (g(0.7) = 2.4725, g(-0.7) = 1.7725).
\begin{tikzpicture}
\fill[acc, opacity=0.10] (-2.4,1.5) rectangle (2.4,2.5);
\fill[siv, opacity=0.18] (-0.7,0) rectangle (0.7,4.2);
\draw[acc, dashed, line width=0.45pt] (-2.4,1.5) -- (2.4,1.5);
\draw[acc, dashed, line width=0.45pt] (-2.4,2.5) -- (2.4,2.5);
\draw[aux] (-0.7,0) -- (-0.7,4.2);
\draw[aux] (0.7,0) -- (0.7,4.2);
\draw[siv, densely dotted] (-2.4,2) -- (2.4,2);
\draw[->] (-2.6,0) -- (2.7,0) node[right]{$h$};
\draw[->] (0,0) -- (0,4.4) node[above]{$g$};
\draw[domain=-2:2, smooth, samples=60] plot (\x, {2+0.5*\x+0.25*\x*\x});
\draw[acc, line width=1.4pt, domain=-0.7:0.7, smooth, samples=20] plot (\x, {2+0.5*\x+0.25*\x*\x});
\node[hole] at (0,2) {};
\node[left] at (-2.4,2.5) {$d + \eta$};
\node[left] at (-2.4,2) {$d$};
\node[left] at (-2.4,1.5) {$d - \eta$};
\node[below] at (-0.7,0) {$-\delta$};
\node[below] at (0.7,0) {$\delta$};
\node[below] at (0,0) {$0$};
\end{tikzpicture}
